\section{Method}
\label{sec:method}
Section~\ref{sec:analysis} motivates a policy representation that preserves the task conditions needed for expert bimanual action selection.
Our framework implements this principle through three connected components, as summarized in Fig.~\ref{fig:architecture}.
First, a task-relation policy transformer combines RGB observations, semantic structure, proprioception, and supervised task queries to predict action chunks.
Second, quality-aware iterative semantic supervision expands a small reviewed seed set without repeatedly invoking a video foundation model.
Third, mask-derived task descriptors supervise the query representation according to the reliability of the semantic evidence from which each descriptor is computed.

\subsection{Task-Relation Policy Transformer}
\label{sec:methodpolicy}
The policy assigns complementary roles to semantic inputs and task-query supervision.
Semantic images expose object-level structure and changes in the deformable occluder to the visual encoder.
Task queries directly supervise selected decision variables in the representation used for action generation.
RGB and proprioception retain appearance, contact, deformation, and robot-configuration cues not captured by category masks.

\subsubsection{Joint RGB-Semantic Encoding}
For each camera view $v$, a frozen segmenter maps the RGB image $I_t^v$ to a class map containing the occluder, target object, goal region, tool, and two robot arms.
Class-dependent component filtering removes small or excess connected components, after which a fixed color palette converts the class map into a three-channel semantic image $S_t^v$.
Policy training uses offline semantic masks from the expanded semantic dataset, while deployment applies the frozen segmenters to the current RGB observations.

Each RGB and semantic stream is processed independently by the shared ACT visual backbone~\cite{zhao2023act}.
The resulting visual tokens are combined with the two-arm proprioceptive state and a set of learned task-query tokens in the transformer encoder.
The ACT decoder attends to the complete encoder memory and predicts a coordinated action chunk.
Thus, the semantic stream exposes changes such as target revelation, occluder displacement, and tool motion without replacing the appearance information in RGB.

\subsubsection{Shared Task Queries and View-Specific Heads}
We use one learned query for each of three task relations: target visibility, tool--target displacement, and target--goal-region distance.
Let $z_t^{(j)}$ denote the encoder output of query $j$.
The same query representation is read by independent heads for the front and side views,
\begin{equation}
 \widehat G_t^{(j,v)}=D_{j,v}\!\left(z_t^{(j)}\right).
 \label{eq:viewqueryheads}
\end{equation}
The shared query integrates complementary evidence across views, while view-specific heads preserve quantities defined in different image coordinates.
The encoded query tokens remain in the encoder memory attended by the action decoder.
Descriptor predictions provide auxiliary supervision and are not appended to the policy input.

\subsection{Quality-Aware Iterative Semantic Supervision}
\label{sec:methodseg}
The semantic images and descriptor targets both depend on mask quality.
We therefore expand the mask annotations through cumulative self-training while screening transient segmentation failures at the video level.

\subsubsection{Reviewed Seeds and Cumulative Expansion}
We first choose a small set of videos that spans key interaction states, including initial occlusion, target exposure, tool contact, placement, and covering restoration.
Bounding boxes and positive and negative point prompts are placed on event-centered keyframes.
Bidirectional Segment Anything Model 2 (SAM2) propagation~\cite{ravi2024sam2}, forward--backward consistency checks, and manual correction produce the reviewed seed mask videos.
The remaining training episodes are task-stratified, randomized, and divided into three disjoint batches.
Each round after the first introduces approximately twice as many new episodes as the preceding round.
The seed episodes and a frozen episode-level benchmark are excluded from the pseudo-label pool.
Algorithm~\ref{alg:iterative-semantic-annotation} summarizes the cumulative expansion.
The benchmark tracks cross-generation segmentation performance.

\begin{figure}[t]
\refstepcounter{troalgorithm}
\label{alg:iterative-semantic-annotation}
\small
\hrule
\vspace{2pt}
\noindent\textbf{Algorithm \thetroalgorithm: Quality-Aware Cumulative Annotation}
\vspace{2pt}
\hrule
\vspace{2pt}
\begin{tabular}{@{}l@{\ }>{\raggedright\arraybackslash}p{0.82\columnwidth}@{}}
\textbf{Input:} & seed masks $\mathcal S_0$, video batches $\mathcal B_{1:K}$, and benchmark $\mathcal E_{\mathrm{seg}}$\\
\textbf{Output:} & segmenter $f_K$ and quality records $\mathcal Q$\\
\end{tabular}
\begin{algorithmic}
\STATE $f_0\leftarrow\operatorname{Train}_{\mathrm{ImageNet}}(\mathcal S_0)$, $\mathcal S\leftarrow\varnothing$, $\mathcal Q\leftarrow\varnothing$
\FOR{$k=1,\ldots,K$}
 \STATE $\widetilde{\mathcal M}_k\leftarrow\operatorname{Predict}(f_{k-1},\mathcal B_k)$
 \STATE $(\bm\upsilon_k,\bm q_k)\leftarrow\operatorname{Quality}(\widetilde{\mathcal M}_k)$
 \STATE $\mathcal A_k\leftarrow\{t:\exists c,v,\alpha_{t,c,v}=1\}$
 \STATE $\mathcal R_k\leftarrow\operatorname{Dilate}\!\left(
 \mathcal A_k\cup\operatorname{GateFail}(\widetilde{\mathcal M}_k)\right)$
 \FOR{$t\in\mathcal B_k$}
  \IF{$t\notin\mathcal R_k$}
   \STATE $\mathcal S\leftarrow\mathcal S\cup\{\widetilde{\mathcal M}_{k,t}\}$, $\mathcal Q\leftarrow\mathcal Q\cup\{\bm q_{k,t}\}$
  \ENDIF
 \ENDFOR
 \STATE $f_k\leftarrow\operatorname{Train}_{\mathrm{ImageNet}}(\mathcal S)$
 \STATE $\operatorname{Evaluate}(f_k,\mathcal E_{\mathrm{seg}})$
\ENDFOR
\end{algorithmic}
\hrule
\end{figure}
Here, $\operatorname{Train}_{\mathrm{ImageNet}}$ denotes training a per-view ResNet U-Net initialized with ImageNet-pretrained encoder weights.
$\operatorname{GateFail}$ denotes excessive structural correction or insufficient foreground, while $\operatorname{Dilate}$ extends each rejected frame to a short temporal neighborhood.

\subsubsection{Video-Level Quality Screening}
Framewise confidence alone conflates genuine task transitions with transient segmentation errors.
For each semantic class $c$ and view $v$, we therefore analyze complete-video trajectories using four evidence groups:
mask area, spatial shape from centroid and bounding-box changes, topology from connected components and holes, and motion consistency from flow-aligned mask overlap.
Area reveals sudden disappearance or expansion.
Centroid and bounding-box trajectories expose implausible location or extent jumps.
Connected components, hole counts, and hole area reveal fragmentation or merging that area alone can miss.
For motion, we compute $I_{t,c,v}$, the intersection over union between the current mask and the preceding mask warped through forward--backward-consistent optical flow.
A drop below its robust expected value gives the motion deviation $d_{t,c,v}^{\mathrm{motion}}=[\widehat I_{t,c,v}-I_{t,c,v}]_+$.
Missing or unreliable flow supplies no motion evidence.

Each continuous diagnostic is compared with a piecewise robust temporal trend so that a persistent change, such as deliberate occluder displacement or target exposure, establishes a new state.
Let $d_{t,c,v}^{g}$ be the resulting normalized deviation for evidence group $g$, and let $\rho_{t,c,v}^{g}$ be its local median absolute residual scaled by $1.4826$.
The excess over the group-specific tolerance $\epsilon_g$ is
\begin{equation}
 \upsilon_{t,c,v}^{g}
 =\left[
 \frac{d_{t,c,v}^{g}}
 {\max(\epsilon_g,\kappa\rho_{t,c,v}^{g})}-1
 \right]_+,
 \qquad \kappa=4,
 \label{eq:qualityevidence}
\end{equation}
where $[u]_+=\max(0,u)$.
Component- or hole-count mismatches directly contribute positive topology evidence.
The four evidence magnitudes define the class-specific quality score
\begin{equation}
 q_{t,c,v}^{\mathrm{seg}}
 =\exp\!\left[-\min\!\left(\tau,
 \sum_g\upsilon_{t,c,v}^{g}\right)\right],
 \qquad \tau=4.
 \label{eq:segquality}
\end{equation}
The score equals one when no evidence departs from its expected trend and decreases as the combined deviation grows.
The multi-evidence anomaly decision is $\alpha_{t,c,v}=\ind[\sum_g\ind[\upsilon_{t,c,v}^{g}>0]\ge2]$.
The joint gate in Algorithm~\ref{alg:iterative-semantic-annotation} therefore prevents an isolated area or motion change from invalidating a label.
The score provides a relative reliability weight rather than a calibrated probability of correctness.
All quality assessment is performed offline.
At deployment, the final frozen segmenters use only the current images.

\subsection{Task Descriptors and Reliability-Aware Supervision}
\label{sec:methodrelations}
The task queries retain information that is directly useful for deciding whether to continue exposure, approach the target, or complete placement.
We instantiate this information using target visibility and two spatial relations.
The more detailed occluder configuration remains available through the semantic stream rather than being compressed into an additional scalar descriptor.

\subsubsection{Task-Descriptor Construction}
Let $M_{O,t}^v$, $M_{T,t}^v$, and $M_{R,t}^v$ denote the target-object, tool, and goal-region masks in view $v$ at time $t$.
For an image of width $W_I$ and height $H_I$, let $\bar c_{O,t}^v$ be the normalized target centroid, $\bar r_{T,t}^v$ a normalized tool reference point, $\mathcal P_v$ a fixed goal-region polygon estimated from reliable training masks, and $D_I=\sqrt{(W_I-1)^2+(H_I-1)^2}$.
The descriptor targets are
\begin{equation}
 \begin{aligned}
 \widetilde G_{t}^{(\mathrm{vis},v)}
 &=\frac{|M_{O,t}^v|}{W_IH_I},\\
 \widetilde G_{t}^{(\mathrm{tool},v)}
 &=\clip\!\left(\bar c_{O,t}^v-\bar r_{T,t}^v,-1,1\right),\\
 \widetilde G_{t}^{(\mathrm{goal},v)}
 &=\frac{1}{D_I}\max_{\xi\in M_{O,t}^v}
   \operatorname{dist}(\xi,\mathcal P_v).
 \end{aligned}
 \label{eq:methodgeometry}
\end{equation}
The current goal-region mask determines descriptor validity and mask reliability, while the fixed polygon $\mathcal P_v$ provides a stable spatial reference.
The visible-area ratio measures the degree of target exposure.
The signed two-dimensional displacement preserves the direction from the execution end toward the target.
In the front view, the tool reference is the left endpoint of the tool mask's principal-axis segment.
In the side view, where endpoint orientation is not reliably identifiable, the tool centroid is used.
Distance to $\mathcal P_v$ is zero inside the goal region, so the maximum in~\eqref{eq:methodgeometry} becomes zero only when all visible target pixels lie inside it.
This retains completion-relevant information that a centroid distance could miss.

A reliably empty target mask is a valid zero-visibility label.
The two geometric relations are undefined when their required objects or reference geometry are unavailable, so they are masked rather than assigned a zero value.
These descriptors encode observable image-plane task relations, while unobserved three-dimensional contact and deformation lie outside their scope.

\subsubsection{Descriptor Reliability and Joint Objective}
Mask-derived descriptor targets inherit segmentation errors.
We therefore weight each descriptor using the semantic classes required for its computation.
Let $\mathcal C_j$ be those classes and $m_{t,j,v}^{G}$ indicate that the geometry is defined.
We use
\begin{equation}
 \varpi_{t,j,v}^{G}
 =m_{t,j,v}^{G}
  \prod_{c\in\mathcal C_j}(1-\alpha_{t,c,v})
  \min_{c\in\mathcal C_j}q_{t,c,v}^{\mathrm{seg}}.
 \label{eq:relationadmission}
\end{equation}
Visibility uses target quality.
Tool--target displacement uses the lower of target and tool quality, while target--goal distance uses the lower of target and goal-region quality.
The weighting depends only on the classes required by each descriptor.
Occluder changes therefore leave valid spatial labels active, while the semantic input retains the occluder structure.
An anomaly or undefined relation gives zero auxiliary weight without removing the corresponding action demonstration.

For action supervision, let $u_{n,h,d}^{\mathrm E}$ and $\widehat u_{n,h,d}$ denote the dataset-normalized expert and predicted joint-command components for sample $n$, prediction step $h$, and command dimension $d$.
Let $p_{n,h}\in\{0,1\}$ mask padded steps, $H_a$ be the prediction horizon, $d_u$ be the joint-command dimension, and $N_b$ be the minibatch size.
During training, the conditional variational encoder also represents each demonstrated action sequence by a Gaussian with mean $\mu_n$ and diagonal variance $\sigma_n^2$.

With $V$ views, three descriptor groups, and component-averaged Smooth L1 loss $\ell_G$, the descriptor and complete training objectives are
\begin{align}
 \mathcal L_{\mathrm{rel}}
 &=\frac{1}{3V}\sum_{j,v}
   \frac{\sum_n\varpi_{n,j,v}^{G}
      \ell_G(\widehat G_n^{(j,v)},\widetilde G_n^{(j,v)})}
        {\max(1,\sum_n m_{n,j,v}^{G})},\nonumber\\
 \mathcal L
 &=\frac{1}{N_bH_ad_u}
  \sum_{n=1}^{N_b}\sum_{h=1}^{H_a}\sum_{d=1}^{d_u}
  p_{n,h}\left|\widehat u_{n,h,d}-u_{n,h,d}^{\mathrm E}\right|
  \nonumber\\
 &\quad+\frac{\beta_{\mathrm{KL}}}{2N_b}
  \sum_{n=1}^{N_b}\sum_r
  \left(\mu_{n,r}^{2}+\sigma_{n,r}^{2}
  -\log\sigma_{n,r}^{2}-1\right)
  +\lambda_{\mathrm{rel}}\mathcal L_{\mathrm{rel}}.
 \label{eq:methodbaseloss}
\end{align}
The denominator counts valid descriptor targets rather than the sum of their quality weights, so uniformly uncertain labels reduce the auxiliary gradient instead of being renormalized to full strength.
The segmenters remain frozen during policy learning.
Together, the shared queries, view-specific heads, and reliability-aware objective provide direct supervision for preserving the selected task information in the action-generating representation, as motivated by~\eqref{eq:expertreconstruction}.
